\documentclass[border=4pt]{standalone}
\usepackage[siunitx]{circuitikz}

\begin{document}
\begin{circuitikz}[european resistors]
  \draw (-2,0) node[ground] {}
    to[V, l_={$V_S=\SI{10.0}{\volt}$}] (-2,4)
    -- (5.0,4);
  \node[font=\small] at (-2.35,3.35) {$+$};
  \node[font=\small] at (-2.35,0.65) {$-$};

  \draw (1.0,4)
    to[R, l_={$R_1\;\SI{1.00}{\kilo\ohm}$}] (1.0,2.35)
    to[R, l_={$R_2\;\SI{2.00}{\kilo\ohm}$}] (1.0,0)
    -- (-2,0);

  \draw (5.0,4)
    to[R, l={$R_3\;\SI{1.50}{\kilo\ohm}$}] (5.0,2.35)
    to[R, l={$R_4\;\SI{3.30}{\kilo\ohm}$}] (5.0,0)
    -- (1.0,0);

  \draw (1.0,2.35)
    to[R, l={$R_5\;\SI{4.70}{\kilo\ohm}$},
       i>^=$I_{12}$] (5.0,2.35);

  % Explicit dots at every junction of three or more conductors.
  \node[circ] at (1.0,4) {};
  \node[circ] at (5.0,4) {};
  \node[circ] at (1.0,2.35) {};
  \node[circ] at (5.0,2.35) {};
  \node[circ] at (1.0,0) {};
  \node[circ] at (5.0,0) {};

  % Test points are led away from components and conductors.
  \draw (1.0,2.35) -- (0.15,2.35)
    node[ocirc, label={[font=\small]left:{TP1: $v_1$}}] {};
  \draw (5.0,2.35) -- (5.85,2.35)
    node[ocirc, label={[font=\small]right:{TP2: $v_2$}}] {};
  \draw (5.0,0) -- (5.85,0)
    node[ocirc, label={[font=\small]right:{TP0: $0$~V}}] {};
  \draw (5.0,4) -- (5.85,4)
    node[ocirc, label={[font=\small]right:{TP$_S$: $V_S$}}] {};

  \node[font=\small] at (3.0,-0.35) {reference node};
\end{circuitikz}
\end{document}
